\section{Recurrences and Master Theorem}
Given a \textcolor{Blue}{recursive algorithm} $\mathcal{A}$, we can express its
running time as a \textcolor{Bittersweet}{\textbf{recurrence relation}} $T(n)$,
describing its running time in terms of the that of its recursive calls.

\subsection{Telescoping}
\textcolor{Blue}{Telescoping} involves \textcolor{Bittersweet}{unrolling} the
recurrence until we reach the base case, and then summing up the contributions
from each level of recursion.
\[\textcolor{ForestGreen}{\text{e.g.}\quad T\left( n \right)
  =\begin{cases}\Theta\left( 1 \right) &\text{if }n\le 1\\2T\left( \frac{n}{2}
\right) +n&\text{if }n>1\end{cases}}.\] 
\begin{multicols}{2}
\begin{align*}
  \frac{T(n)}{n} =& \frac{T\left( \frac{n}{2} \right)}{\frac{n}{2}} + 1\\
  \frac{T\left( \frac{n}{2} \right)}{\frac{n}{2}} =& \frac{T\left( \frac{n}{4} \right)}{\frac{n}{4}} + 1\\
       \vdots\\
  \frac{T\left( 2 \right) }{2} =& \frac{T\left( 1 \right)}{1} + 1
\end{align*}
\vspace{12mm}
\begin{align*}
  \implies \frac{T(n)}{n} =& \frac{T\left( 1 \right)}{1} + \log n\\
  =& O(\log n)
\end{align*}
\end{multicols}
\vspace{-4mm}
\subsection{Substitution}
\textcolor{Blue}{Substitution} requires \textcolor{Bittersweet}{guessing} the
form of the solution, and then \textcolor{Bittersweet}{proving} it by
induction.
\[\textcolor{ForestGreen}{\text{e.g.}\quad T\left( n \right) =\begin{cases}c&\text{if }n\le 1\\4T\left( \lfloor\frac{n}{2}\rfloor \right) +n&\text{if }n>1\end{cases}}.\] 
\begin{align*}
  \text{Induction Hypothesis: } T(k) \le& (c+1)k^2-n\text{ for all }k<n
\end{align*}
\begin{align*}
  \text{Induction Step: } T(n) =& 4T\left( \lfloor\frac{n}{2}\rfloor \right) +n\\
  \le& 4(c+1)\left( \frac{n}{2} \right)^2 -4\left( \frac{n}{2} \right) + n\\
  =& (c+1)n^2 - n
\end{align*}
\paragraph{Making good guesses}
Consider $T(n) =$ \textcolor{ForestGreen}{\textbf{Recursive calls}} $+$
\textcolor{Bittersweet}{\textbf{``Merging'' time}}. Substituting the guess
back into the recurrence, and compare to the guess each of the contributions:
\begin{enumerate}
  \item Time contributed by \textcolor{ForestGreen}{recursive calls} (Including
    constants.)
  \item Time contributed by \textcolor{Bittersweet}{``merging'' time}
    (Asymptotically.)
\end{enumerate}
Comparing these contributions to the guess:
\begin{itemize}
  \item If either is larger than the guess, try a \textbf{larger guess}.
  \item If both are smaller than the guess, try a \textbf{smaller guess}.
  \item If both are equal to the guess, add a $\log{n}$ factor.
\end{itemize}

\subsection{Recursion Tree}
\textcolor{Blue}{Recursion tree} involves \textcolor{Bittersweet}{visualising}
the recursive calls as a tree, and then summing up the contributions from each
level of recursion.
\[\textcolor{ForestGreen}{\text{e.g.}\quad T\left( n \right)
  =\begin{cases}\Theta\left( 1 \right) &\text{if }n\le 1\\2T\left(\frac{n}{2}
\right) +cn&\text{if }n>1\end{cases}}.\] 
\begin{center}
  \includegraphics[width=0.24\textwidth]{recursion_tree.png}
\end{center}
\vspace{-4mm}

\subsection{Master Theorem}
\textcolor{Blue}{Master Theorem} provides a \textcolor{Bittersweet}{closed-form solution} to recurrences of the form
\[T(n)=aT\left( \frac{n}{b} \right) +f(n)\] 
where $a\ge 1$ and $b>1$ are constants, and $f(n)\in\Omega\left( 1 \right) $.
For \textcolor{Blue}{$d=\log_b{a}$}, and  $\epsilon>0$:
\begin{itemize}
  \item $f(n)\in O\left( n^{d-\epsilon} \right) \implies T(n)\in \Theta\left(
    n^d \right) $.
  \item $f(n)\in \Theta\left( n^d \log^{k}{n} \right) \text{ for some
    }k\ge 0\implies T(n)\in \Theta\left( n^d\log^{k+1} n \right) $.
  \item $f(n)\in \Omega\left( n^{d+\epsilon} \right) $, and $af\left(
    \frac{n}{b} \right) \le cf(n)$ for some constant $c<1\implies T(n)\in
    \Theta(f(n))$.
\end{itemize}
Floors and ceilings can be ignored in the Master Theorem as they
\textcolor{Bittersweet}{do not affect the asymptotic behaviour} of the
recurrence.
\paragraph{Extended Master Theorem}
Consider a recurrence of the form
\[T(n)=aT\left( \frac{n}{b} \right) +f(n)\] 
where $a\ge 1$ and $b>1$ are constants. For $d=\log_b{a}$,
and \textcolor{Blue}{$f(n)\in\Theta\left(n^{d}\log^{k}{n} \right) $}:
\begin{itemize}
  \item If $k>-1$, then $T(n)\in \Theta\left( n^d\log^{k+1} n \right) $.
  \item If $k=-1$, then $T(n)\in \Theta\left (n^d\log\log n) \right) $.
  \item If $k<-1$, then $T(n)\in \Theta\left( n^d \right) $.
\end{itemize}
\textbf{Proof:} Consider $T(n)=aT\left( \frac{n}{b} \right) +n^{d}\log^{k}{n}$.
Dividing both sides by $n^d$:
\begin{align*}
  \frac{T(n)}{n^{d}} =& \frac{T\left( \frac{n}{b} \right)}{{\left( \frac{n}{b}
  \right)}^{d}} + \log^{k}{n}\\
        \frac{T\left( \frac{n}{b} \right)}{{\left( \frac{n}{b} \right)}^{d} } =& \frac{T\left( \frac{n}{b^{2}} \right)}{{\left( \frac{n}{b^{2}} \right)}^{d} } + \log^{k}{\left( \frac{n}{b}\right)} \\
        \frac{T\left( \frac{n}{b^{2}} \right)}{{\left( \frac{n}{b^{2}} \right)}^{d} } =& \frac{T\left( \frac{n}{b^{3}} \right)}{{\left( \frac{n}{b^{3}} \right)}^{d} } + \log^{k}{\left( \frac{n}{b^{2}}\right)} \\
       \vdots
\end{align*}
\[\frac{T\left( n \right) }{n^{d}}=\log^{k}{n}+\log^{k}{\left( \frac{n}{b} \right) }+\log^{k}{\left( \frac{n}{b^{2}} \right) }+\cdots+\Theta\left( 1 \right) .\] 
Assuming logarithms are base $b$:
\begin{align*}
\frac{T\left( n \right) }{n^{d}}={\left( \log{n} \right)}^{k} &+{\left( \log{ n  }-1 \right)}^{k} +{\left( \log{ n  }-2 \right)}^{k}\\&+\cdots+\Theta\left( 1 \right)
\end{align*}
If $k<-1$, we have the p-series test with $p=-k>1$, and thus:
\[\sum_{x=1}^{\log{n}} x^{k}=\Theta\left( 1 \right).\]
If $k=-1$, we have the harmonic series, and thus:
\[\sum_{x=1}^{\log{n}} x^{-1}=\Theta\left( \log\log{n} \right).\]
If $k>-1$, we have (sum of first $n$ integers to the $k$-th power):
\begin{align*}
  \sum_{x=1}^{\log{n}} x^{k}&=\frac{\log^{k+1}{n}}{k+1}+O\left( \log^{k}{n} \right)\\
                            &=\Theta\left( \log^{k+1}{n} \right)
\end{align*}
